\documentclass[12pt, a4paper]{article}
% --- 1. Cấu hình Tiếng Việt và Font chữ chuẩn ---
\usepackage[utf8]{inputenc}
\usepackage[T5]{fontenc}
\usepackage[vietnamese]{babel}
\usepackage{mathptmx} % Font Times New Roman đồng bộ cả chữ và công thức toán, không gây xung đột gói

\makeatletter
\renewcommand{\normalsize}{\@setfontsize\normalsize{13pt}{16pt}}
\makeatother
\normalsize

% --- 2. Gói Toán học & Ký hiệu bổ trợ ---
\usepackage{amsmath, amssymb, amsfonts, mathrsfs, amsthm}
\usepackage{graphicx}
\usepackage{fontawesome}
\usepackage{booktabs, tabularx, array, multirow}
\usepackage{enumitem}

% --- 3. Đồ họa TikZ, Bảng biến thiên & Hình không gian ---
\usepackage{tikz}
\usetikzlibrary{calc, intersections, angles, quotes, patterns, arrows.meta, 3d}
\usepackage{pgfplots}
\pgfplotsset{compat=1.18}
\usepackage{tkz-euclide}
\usepackage{tkz-tab}

\renewcommand{\vec}[1]{\overrightarrow{#1}}

% --- 4. Hộp sư phạm (Tcolorbox) ---
\usepackage[many]{tcolorbox}
\tcbuselibrary{skins, breakable}

% --- 5. Căn lề chuẩn văn bản giáo án ---
\usepackage{geometry}
\geometry{
    a4paper,
    left=25mm,
    right=15mm,
    top=20mm,
    bottom=20mm
}

% --- 6. Header / Footer chuẩn hóa ---
\usepackage{fancyhdr}
\usepackage{lastpage}
\pagestyle{fancy}
\fancyhf{}

\fancyhead[L]{\small \textit{Kế hoạch bài dạy Toán 11}}
\fancyhead[R]{\textbf{Trang \thepage/\pageref{LastPage}}}
\fancyfoot[L]{\small \textit{Tổ Toán - Tin}}
\fancyfoot[R]{\small \textit{Trường THCS\&THPT Hồng Vân}}
\renewcommand{\headrulewidth}{0.4pt}
\renewcommand{\footrulewidth}{0.4pt}

% --- 7. Giãn dòng & Giãn đoạn theo chuẩn Nghị định 30/2020/NĐ-CP ---
\usepackage{setspace}
\setstretch{1.15}
\setlength{\parindent}{1.0cm}
\setlength{\parskip}{5pt}

% Định nghĩa hộp sư phạm chuyên dụng
\newtcolorbox{pedabox}[2][]{
    enhanced,
    breakable,
    colback=blue!3!white,
    colframe=blue!65!black,
    fonttitle=\bfseries\small,
    title={#2},
    arc=2mm,
    boxrule=0.6pt,
    left=3mm, right=3mm, top=2mm, bottom=2mm,
    #1
}

\newtcolorbox{aibox}[2][]{
    enhanced,
    breakable,
    colback=teal!4!white,
    colframe=teal!70!black,
    fonttitle=\bfseries\small,
    title={#2},
    arc=2mm,
    boxrule=0.6pt,
    left=3mm, right=3mm, top=2mm, bottom=2mm,
    #1
}

\newtcolorbox{scaffoldbox}[2][]{
    enhanced,
    breakable,
    colback=orange!4!white,
    colframe=orange!80!black,
    fonttitle=\bfseries\small,
    title={#2},
    arc=2mm,
    boxrule=0.6pt,
    left=3mm, right=3mm, top=2mm, bottom=2mm,
    #1
}

\newtcolorbox{stembox}[2][]{
    enhanced,
    breakable,
    colback=purple!4!white,
    colframe=purple!75!black,
    fonttitle=\bfseries\small,
    title={#2},
    arc=2mm,
    boxrule=0.6pt,
    left=3mm, right=3mm, top=2mm, bottom=2mm,
    #1
}

\begin{document}

\begin{center}
    {\fontsize{14pt}{17pt}\selectfont \textbf{KẾ HOẠCH BÀI DẠY MÔN TOÁN LỚP 11}}\\[6pt]
    {\fontsize{16pt}{19pt}\selectfont \textbf{BÀI 18. LUỸ THỪA VỚI SỐ MŨ THỰC}}\\[4pt]
    {\fontsize{12pt}{15pt}\selectfont \textbf{Môn học: Toán 11} \ \textbar \ \textbf{Thời lượng: 2 tiết (Tiết 55, 56 theo PPCT | Tuần 19)}}
\end{center}

\vspace{5pt}

\section*{I. MỤC TIÊU}

\subsection*{1. Kiến thức, Năng lực}

\begin{itemize}[leftmargin=1.2cm]
    \item \textbf{Yêu cầu cần đạt (Nguyên văn CT GDPT 2018 Toán 11):}
    \begin{itemize}[leftmargin=0.6cm]
        \item Nhận biết được khái niệm luỹ thừa với số mũ nguyên của một số thực khác $0$; luỹ thừa với số mũ hữu tỉ và luỹ thừa với số mũ thực của một số thực dương.
        \item Giải thích được các tính chất của phép tính luỹ thừa với số mũ nguyên, luỹ thừa với số mũ hữu tỉ và luỹ thừa với số mũ thực.
        \item Sử dụng được tính chất của phép tính luỹ thừa trong tính toán các biểu thức số và rút gọn các biểu thức chứa biến.
        \item Tính được giá trị biểu thức số có chứa phép tính luỹ thừa bằng cách sử dụng máy tính cầm tay.
        \item Giải quyết được một số vấn đề có liên quan đến môn học khác hoặc thực tiễn gắn với phép tính luỹ thừa (ví dụ: bài toán tăng trưởng vi khuẩn, lãi kép, phân rã phóng xạ).
    \end{itemize}

    \item \textbf{Năng lực Toán học đặc thù:}
    \begin{itemize}[leftmargin=0.6cm]
        \item \textit{Năng lực tư duy và lập luận toán học:} So sánh, phân biệt sự mở rộng phạm vi số mũ từ số tự nhiên $\to$ số nguyên $\to$ số hữu tỉ $\to$ số thực kèm theo sự thu hẹp tương ứng của tập xác định cơ số; lập luận chặt chẽ khi biến đổi, so sánh các biểu thức lũy thừa cùng cơ số hoặc khác cơ số.
        \item \textit{Năng lực mô hình hóa toán học:} Thiết lập mô hình tăng trưởng hàm mũ $N(t) = N_0 \cdot 2^{t/T}$ hoặc mô hình lãi kép $A = P(1 + r)^n$ từ các bài toán thực tiễn trong sinh học, kinh tế tài chính.
        \item \textit{Năng lực giải quyết vấn đề toán học:} Phân tích cấu trúc biểu thức chứa căn và lũy thừa nhiều tầng để đưa về cùng một cơ số; lựa chọn thứ tự biến đổi hợp lý để triệt tiêu nhân tử chung và rút gọn biểu thức chứa biến.
        \item \textit{Năng lực giao tiếp toán học:} Sử dụng chính xác các thuật ngữ toán học: cơ số, số mũ, căn bậc $n$, lũy thừa với số mũ hữu tỉ, vô tỉ; trình bày các bước chứng minh và biến đổi biểu thức toán học rõ ràng, mạch lạc.
        \item \textit{Năng lực sử dụng công cụ, phương tiện học toán:} Sử dụng thành thạo máy tính cầm tay (MTCT) Casio fx-580VN X: phím lũy thừa \fbox{$x^\blacksquare$}, phím căn bậc $n$ \fbox{$\sqrt[\blacksquare]{\blacksquare}$}, phím số thực $\pi, e$ để tính giá trị biểu thức và kiểm tra kết quả rút gọn bằng phương pháp gán giá trị hoặc lệnh \texttt{CALC}.
    \end{itemize}

    \item \textbf{Năng lực số (Theo Thông tư 02/2025/TT-BGDĐT):}
    \begin{itemize}[leftmargin=0.6cm]
        \item \textit{Mã 5.1NC1a (Ứng dụng công nghệ tính toán số chuyên dụng):} Học sinh sử dụng máy tính cầm tay (MTCT Casio fx-580VN X) để thực hiện tính toán giá trị biểu thức lũy thừa với số mũ hữu tỉ dạng $a^{\frac{m}{n}} = \sqrt[n]{a^m}$, xấp xỉ số mũ vô tỉ qua dãy số hữu tỉ; phân tích và xử lý lỗi toán học (Math ERROR) khi nhập cơ số âm với số mũ không nguyên.
    \end{itemize}

    \item \textbf{Năng lực AI (Theo Quyết định 2422/QĐ-BGDĐT):}
    \begin{itemize}[leftmargin=0.6cm]
        \item \textit{Mã 11.C5.1 (Hiểu biết nền tảng toán học của công nghệ AI):} Học sinh nêu được vai trò then chốt của phép toán lũy thừa và hàm mũ trong việc biểu diễn độ phức tạp thuật toán, tốc độ xử lý lũy thừa ma trận trọng số và hàm kích hoạt phi tuyến tính (Softmax, Sigmoid) trong mạng nơ-ron nhân tạo (Artificial Neural Networks).
    \end{itemize}

    \item \textbf{Năng lực chung:}
    \begin{itemize}[leftmargin=0.6cm]
        \item \textit{Tự chủ và tự học:} Chủ động hệ thống hóa bảng tính chất của lũy thừa; tự giác phát hiện lỗi sai thường gặp về điều kiện cơ số.
        \item \textit{Giao tiếp và hợp tác:} Tương tác tích cực trong thảo luận nhóm, giải thích quy tắc đưa căn thức về lũy thừa số mũ hữu tỉ cho bạn học.
    \end{itemize}
\end{itemize}

\subsection*{2. Phẩm chất}

\begin{itemize}[leftmargin=1.2cm]
    \item \textbf{Chăm chỉ:} Kiên trì rèn luyện kỹ năng biến đổi lũy thừa, không ngại các biểu thức phức tạp nhiều tầng căn.
    \item \textbf{Trung thực:} Tự lực thực hiện các bài toán rút gọn, sử dụng máy tính cầm tay để kiểm chứng kết quả chứ không sao chép máy móc.
    \item \textbf{Trách nhiệm:} Có trách nhiệm giải thích rõ ràng căn cứ toán học cho từng bước biến đổi của nhóm mình; tôn trọng kỷ luật phòng học.
\end{itemize}

\vspace{5pt}

\section*{II. THIẾT BỊ DẠY HỌC VÀ HỌC LIỆU}

\begin{itemize}[leftmargin=1.2cm]
    \item \textbf{Giáo viên:}
    \begin{itemize}[leftmargin=0.6cm]
        \item Kế hoạch bài dạy chi tiết 2 tiết theo chuẩn Công văn 5512/BGDĐT-GDTrH.
        \item Hệ thống Slide bài giảng tương tác trực quan, mô phỏng phần mềm mô phỏng MTCT Casio fx-580VN X trên máy tính.
        \item Hệ thống Phiếu học tập phân hóa bậc thang: Phiếu học tập số 1 (Tiết 1) và Phiếu học tập số 2 (Tiết 2).
        \item Bảng Rubric đánh giá năng lực học sinh trong bài học.
    \end{itemize}
    \item \textbf{Học sinh:}
    \begin{itemize}[leftmargin=0.6cm]
        \item Sách giáo khoa Toán 11, vở ghi chép bài học, giấy nháp.
        \item Máy tính cầm tay Casio fx-580VN X (hoặc tương đương).
        \item Đọc trước bài học mới ở nhà theo hướng dẫn của giáo viên.
    \end{itemize}
\end{itemize}

\vspace{5pt}

\section*{III. TIẾN TRÌNH DẠY HỌC (2 TIẾT | TIẾT 55, 56 THEO PPCT)}

\subsection*{\faClockO \ TIẾT 1 (TIẾT 55 THEO PPCT): LŨY THỪA VỚI SỐ MŨ NGUYÊN, CĂN BẬC $n$ VÀ LŨY THỪA VỚI SỐ MŨ HỮU TỈ}

\subsubsection*{1. Hoạt động 1: Mở đầu (Khởi động) (5 phút)}

\noindent \textbf{a) Mục tiêu:}
Tạo mâu thuẫn nhận thức về sự cần thiết mở rộng khái niệm lũy thừa từ số mũ nguyên dương sang số mũ $0$, nguyên âm và hữu tỉ thông qua tình huống thực tế.

\noindent \textbf{b) Nội dung:}
Giáo viên đặt tình huống:
\begin{pedabox}{Tình huống khởi động: Sự phân chia của tế bào vi khuẩn}
Một quần thể vi khuẩn ban đầu có $1000$ cá thể, sau mỗi giờ số lượng vi khuẩn tăng gấp đôi.
\begin{enumerate}[leftmargin=0.6cm]
    \item Số lượng vi khuẩn sau $1$ giờ, $2$ giờ, $n$ giờ ($n \in \mathbb{N}^*$) là bao nhiêu?
    \item Sau $0$ giờ (thời điểm ban đầu), số lượng vi khuẩn là $1000$. Biểu thức lũy thừa tương ứng là $1000 \cdot 2^0$. Vậy $2^0$ bằng bao nhiêu?
    \item Trước thời điểm ban đầu $1$ giờ (tức thời điểm $t = -1$ giờ), số lượng vi khuẩn là $500 = 1000 \cdot \dfrac{1}{2} = 1000 \cdot 2^{-1}$. Vậy $2^{-1}$ bằng bao nhiêu?
    \item Sau nửa giờ ($t = \dfrac{1}{2}$ giờ), số lượng vi khuẩn tăng bao nhiêu lần? Đại lượng $2^{\frac{1}{2}}$ có ý nghĩa toán học như thế nào?
\end{enumerate}
\end{pedabox}

\noindent \textbf{c) Sản phẩm:}
Câu trả lời của học sinh:
\begin{itemize}
    \item Sau $n$ giờ, số vi khuẩn là $1000 \cdot 2^n$.
    \item Thời điểm ban đầu $t = 0$: $1000 \cdot 2^0 = 1000 \implies 2^0 = 1$.
    \item Thời điểm $t = -1$: $1000 \cdot 2^{-1} = 500 \implies 2^{-1} = \dfrac{1}{2}$.
    \item Thời điểm $t = \dfrac{1}{2}$: Học sinh nhận thấy nếu sau hai lần nửa giờ vi khuẩn nhân đôi ($2^{\frac{1}{2}} \cdot 2^{\frac{1}{2}} = 2^1 = 2$), suy ra $2^{\frac{1}{2}} = \sqrt{2}$. Xuất hiện nhu cầu định nghĩa lũy thừa với số mũ hữu tỉ!
\end{itemize}

\noindent \textbf{d) Tổ chức thực hiện:}
\begin{itemize}[leftmargin=0.6cm]
    \item \textbf{Giao nhiệm vụ:} Giáo viên trình chiếu tình huống, yêu cầu học sinh làm việc cá nhân trong 2 phút rồi trao đổi cặp đôi.
    \item \textbf{Thực hiện nhiệm vụ:} Học sinh tính toán, quan sát quy luật giảm dần của số mũ: $2^3 = 8, 2^2 = 4, 2^1 = 2, 2^0 = 1, 2^{-1} = \dfrac{1}{2}, \dots$
    \item \textbf{Báo cáo, thảo luận:} Đại diện 1 học sinh trả lời, các học sinh khác nhận xét, bổ sung.
    \item \textbf{Kết luận, nhận định:} Giáo viên dẫn dắt: Để phục vụ mô tả các quá trình biến thiên liên tục trong tự nhiên, toán học cần mở rộng khái niệm lũy thừa sang số mũ nguyên âm, hữu tỉ và số thực.
\end{itemize}

\subsubsection*{2. Hoạt động 2: Hình thành kiến thức (25 phút)}

\paragraph{2.1. Đơn vị kiến thức 1: Lũy thừa với số mũ nguyên (8 phút)}

\noindent \textbf{a) Mục tiêu:}
Nhận biết và giải thích được định nghĩa lũy thừa với số mũ nguyên của một số thực khác $0$; nắm vững điều kiện của cơ số.

\noindent \textbf{b) Nội dung:}
Học sinh tiếp cận định nghĩa lũy thừa với số mũ nguyên qua bảng phân loại điều kiện cơ số:
\begin{pedabox}{Định nghĩa Lũy thừa với số mũ nguyên}
Cho $n$ là một số nguyên dương, $a$ là một số thực.
\begin{enumerate}[leftmargin=0.6cm]
    \item Với $a \in \mathbb{R}$: $a^n = \underbrace{a \cdot a \cdots a}_{n \text{ thừa số } a}$.
    \item Với $a \neq 0$:
    \begin{itemize}
        \item $a^0 = 1$.
        \item $a^{-n} = \dfrac{1}{a^n}$.
    \end{itemize}
\end{enumerate}
\end{pedabox}

\begin{scaffoldbox}{Biện pháp sư phạm: Lưu ý cốt tử về điều kiện cơ số (Scaffolding)}
\textbf{Khắc sâu cho học sinh trung bình - yếu:}
\begin{itemize}
    \item Biểu thức $0^0$ và $0^{-n}$ (với $n \in \mathbb{N}^*$) là \textbf{KHÔNG CÓ NGHĨA}!
    \item Ví dụ: $0^{-2} = \dfrac{1}{0^2} = \dfrac{1}{0}$ (phép chia cho $0$ không xác định).
    \item Vì vậy, khi gặp số mũ nguyên âm hoặc số mũ $0$, điều kiện bắt buộc là \textbf{cơ số phải khác 0} ($a \neq 0$).
\end{itemize}
\end{scaffoldbox}

\noindent \textbf{c) Sản phẩm:}
Học sinh tính nhanh các ví dụ:
\begin{itemize}
    \item $(-3)^0 = 1$; $\left(\dfrac{2}{3}\right)^{-2} = \dfrac{1}{\left(\frac{2}{3}\right)^2} = \dfrac{1}{\frac{4}{9}} = \dfrac{9}{4}$.
    \item $(-2)^{-3} = \dfrac{1}{(-2)^3} = -\dfrac{1}{8}$.
\end{itemize}

\noindent \textbf{d) Tổ chức thực hiện:}
\begin{itemize}[leftmargin=0.6cm]
    \item \textbf{Giao nhiệm vụ:} Giáo viên trình bày định nghĩa, nhấn mạnh điều kiện $a \neq 0$ khi số mũ nguyên không dương. Yêu cầu học sinh tính giá trị các biểu thức trên.
    \item \textbf{Thực hiện nhiệm vụ:} Học sinh làm bài vào vở; giáo viên theo dõi, hướng dẫn học sinh còn lúng túng với dấu âm ở cơ số và số mũ.
    \item \textbf{Báo cáo, thảo luận:} Gọi 2 học sinh lên bảng trình bày; học sinh lớp nhận xét.
    \item \textbf{Kết luận, nhận định:} Giáo viên chuẩn hóa định nghĩa và các tính chất cơ bản: $a^m \cdot a^n = a^{m+n}$, $\dfrac{a^m}{a^n} = a^{m-n}$, $(a^m)^n = a^{m \cdot n}$, $(ab)^n = a^n b^n$, $\left(\dfrac{a}{b}\right)^n = \dfrac{a^n}{b^n}$.
\end{itemize}

\paragraph{2.2. Đơn vị kiến thức 2: Căn bậc $n$ của một số thực (8 phút)}

\noindent \textbf{a) Mục tiêu:}
Nhận biết khái niệm căn bậc $n$ của một số thực; phân biệt tính chất căn bậc chẵn và căn bậc lẻ.

\noindent \textbf{b) Nội dung:}
\begin{pedabox}{Khái niệm và Tính chất của Căn bậc $n$}
Cho số nguyên dương $n \ge 2$ và số thực $a$. Số thực $b$ được gọi là căn bậc $n$ của số thực $a$ nếu $b^n = a$.
\begin{itemize}[leftmargin=0.6cm]
    \item \textbf{Trường hợp $n$ lẻ ($n = 2k + 1, k \in \mathbb{N}^*$):} Mọi số thực $a$ đều có duy nhất một căn bậc $n$, ký hiệu là $\sqrt[n]{a}$.
    \begin{itemize}
        \item $\sqrt[n]{a} > 0$ khi $a > 0$; $\sqrt[n]{a} < 0$ khi $a < 0$; $\sqrt[n]{0} = 0$.
    \end{itemize}
    \item \textbf{Trường hợp $n$ chẵn ($n = 2k, k \in \mathbb{N}^*$):}
    \begin{itemize}
        \item Với $a < 0$: Không tồn tại căn bậc $n$ của $a$.
        \item Với $a = 0$: Có một căn bậc $n$ duy nhất là $0$ ($\sqrt[n]{0} = 0$).
        \item Với $a > 0$: Có đúng hai căn bậc $n$ là hai số đối nhau: giá trị dương ký hiệu là $\sqrt[n]{a}$, giá trị âm ký hiệu là $-\sqrt[n]{a}$.
    \end{itemize}
\end{itemize}
\end{pedabox}

\noindent \textbf{c) Sản phẩm:}
Học sinh hoàn thành câu hỏi nhanh:
\begin{itemize}
    \item $\sqrt[3]{-27} = -3$ vì $(-3)^3 = -27$.
    \item $\sqrt[4]{16} = 2$ (căn bậc $4$ số học); căn bậc $4$ của $16$ gồm hai giá trị là $2$ và $-2$.
    \item $\sqrt[4]{-16}$ không có nghĩa trong tập số thực $\mathbb{R}$.
\end{itemize}

\noindent \textbf{d) Tổ chức thực hiện:}
\begin{itemize}[leftmargin=0.6cm]
    \item \textbf{Giao nhiệm vụ:} Giáo viên yêu cầu học sinh thảo luận cặp đôi: So sánh sự khác nhau giữa căn bậc 3 của $-8$ và căn bậc 2 của $-4$.
    \item \textbf{Thực hiện nhiệm vụ:} Học sinh đối chiếu định nghĩa, rút ra kết luận về căn bậc chẵn và căn bậc lẻ.
    \item \textbf{Báo cáo, thảo luận:} Học sinh nêu rõ lý do tại sao căn bậc chẵn không lấy được với số âm.
    \item \textbf{Kết luận, nhận định:} Giáo viên nhấn mạnh quy tắc khai căn: $\sqrt[n]{a^n} = a$ khi $n$ lẻ, và $\sqrt[n]{a^n} = |a|$ khi $n$ chẵn.
\end{itemize}

\paragraph{2.3. Đơn vị kiến thức 3: Lũy thừa với số mũ hữu tỉ (9 phút)}

\noindent \textbf{a) Mục tiêu:}
Hiểu rõ định nghĩa lũy thừa với số mũ hữu tỉ của một số thực dương; sử dụng thành thạo MTCT Casio fx-580VN X để tính giá trị (Mã 5.1NC1a).

\noindent \textbf{b) Nội dung:}
\begin{pedabox}{Định nghĩa Lũy thừa với số mũ hữu tỉ}
Cho số thực $a > 0$ và số hữu tỉ $r = \dfrac{m}{n}$, trong đó $m \in \mathbb{Z}, n \in \mathbb{N}, n \ge 2$. Khi đó:
$$a^r = a^{\frac{m}{n}} = \sqrt[n]{a^m}$$
Đặc biệt: $a^{\frac{1}{n}} = \sqrt[n]{a}$.
\end{pedabox}

\begin{scaffoldbox}{Giàn giáo sư phạm: Bắt buộc cơ số dương khi số mũ không nguyên}
\textbf{Sai lầm kinh điển cần triệt tiêu:}
\begin{itemize}
    \item Học sinh thường tự ý viết: $(-8)^{\frac{1}{3}} = \sqrt[3]{-8} = -2 \implies$ \textbf{SAI VỀ MẶT ĐỊNH NGHĨA}!
    \item Giải thích: Nếu chấp nhận cơ số âm cho số mũ hữu tỉ, ta sẽ có mâu thuẫn toán học:
    $$-2 = (-8)^{\frac{1}{3}} = (-8)^{\frac{2}{6}} = \sqrt[6]{(-8)^2} = \sqrt[6]{64} = 2 \implies -2 = 2 \ (\text{VÔ LÝ!})$$
    \item Do đó, quy chuẩn toán học quốc tế và CT GDPT 2018 quy định: \textbf{Lũy thừa với số mũ không nguyên chỉ xác định khi cơ số DƯƠNG ($a > 0$)}.
\end{itemize}
\end{scaffoldbox}

\begin{aibox}{Tích hợp Năng lực số (Mã 5.1NC1a): Thao tác bấm MTCT Casio fx-580VN X}
\begin{itemize}[leftmargin=0.6cm]
    \item Hướng dẫn học sinh tính $8^{\frac{2}{3}}$: Bấm phím \fbox{8} $\to$ \fbox{$x^\blacksquare$} $\to$ \fbox{$\dfrac{\blacksquare}{\blacksquare}$} $\to$ \fbox{2} $\to$ \fbox{$\blacktriangledown$} $\to$ \fbox{3} $\to$ \fbox{$=$}. Kết quả màn hình hiển thị: $4$.
    \item Kiểm tra công thức căn thức: Bấm \fbox{$\text{SHIFT}$} $\to$ \fbox{$\sqrt[\blacksquare]{\blacksquare}$} $\to$ nhập $n=3$, bên trong nhập $8^2 \to$ hiển thị $4$.
    \item Nhập thử biểu thức $(-4)^{\frac{1}{2}}$: Máy tính báo \texttt{Math ERROR}. Học sinh giải thích được nguyên nhân do vi phạm điều kiện cơ số dương.
\end{itemize}
\end{aibox}

\noindent \textbf{c) Sản phẩm:}
Học sinh hoàn thành bảng tính:
\begin{center}
\begin{tabular}{|c|c|c|c|c|}
    \hline
    Biểu thức lũy thừa & $16^{\frac{3}{4}}$ & $27^{-\frac{2}{3}}$ & $\left(\dfrac{1}{32}\right)^{0{,}4}$ & $81^{\frac{1}{4}}$ \\ \hline
    Chuyển về căn thức & $\sqrt[4]{16^3} = (\sqrt[4]{16})^3$ & $\dfrac{1}{\sqrt[3]{27^2}} = \dfrac{1}{(\sqrt[3]{27})^2}$ & $\sqrt[5]{\left(\dfrac{1}{32}\right)^2}$ & $\sqrt[4]{81}$ \\ \hline
    Giá trị tính tay & $2^3 = 8$ & $\dfrac{1}{3^2} = \dfrac{1}{9}$ & $\left(\dfrac{1}{2}\right)^2 = \dfrac{1}{4}$ & $3$ \\ \hline
    Kiểm tra MTCT & $8$ & $0{,}1111\dots = \dfrac{1}{9}$ & $0{,}25 = \dfrac{1}{4}$ & $3$ \\ \hline
\end{tabular}
\end{center}

\noindent \textbf{d) Tổ chức thực hiện:}
\begin{itemize}[leftmargin=0.6cm]
    \item \textbf{Giao nhiệm vụ:} Giáo viên trình bày định nghĩa, hướng dẫn thao tác trên MTCT; phát Phiếu học tập số 1 (phần 1).
    \item \textbf{Thực hiện nhiệm vụ:} Học sinh làm việc cá nhân, tính toán bằng tay rồi dùng MTCT bấm kiểm tra đối chiếu.
    \item \textbf{Báo cáo, thảo luận:} Giáo viên gọi học sinh nêu kết quả và thao tác phím bấm tương ứng.
    \item \textbf{Kết luận, nhận định:} Khắc sâu công thức $a^{\frac{m}{n}} = \sqrt[n]{a^m} = (\sqrt[n]{a})^m$ (với $a > 0$).
\end{itemize}

\subsubsection*{3. Hoạt động 3: Luyện tập (10 phút)}

\noindent \textbf{a) Mục tiêu:}
Củng cố kỹ năng biến đổi, rút gọn biểu thức số và biểu thức chứa biến bằng định nghĩa và tính chất của lũy thừa với số mũ hữu tỉ.

\noindent \textbf{b) Nội dung:}
Thực hiện các bài toán trong Phiếu học tập số 1:
\begin{enumerate}[leftmargin=0.6cm]
    \item \textbf{Bài 1:} Tính giá trị biểu thức: $A = \dfrac{2^3 \cdot 2^{-1} + 5^{-3} \cdot 5^4}{10^{-3} \cdot 10^4 - (0{,}1)^0}$.
    \item \textbf{Bài 2:} Rút gọn biểu thức (với $x > 0, y > 0$): $B = \dfrac{x^{\frac{5}{4}} \cdot y + x \cdot y^{\frac{5}{4}}}{\sqrt[4]{x} + \sqrt[4]{y}}$.
\end{enumerate}

\noindent \textbf{c) Sản phẩm:}
Lời giải chi tiết:
\begin{itemize}
    \item \textbf{Bài 1:}
    \begin{align*}
        A &= \dfrac{2^{3 + (-1)} + 5^{-3 + 4}}{10^{-3 + 4} - 1} = \dfrac{2^2 + 5^1}{10^1 - 1} = \dfrac{4 + 5}{10 - 1} = \dfrac{9}{9} = 1.
    \end{align*}
    \item \textbf{Bài 2:}
    Nhận xét: $x^{\frac{5}{4}} = x^{1 + \frac{1}{4}} = x \cdot x^{\frac{1}{4}} = x \cdot \sqrt[4]{x}$; tương tự $y^{\frac{5}{4}} = y \cdot \sqrt[4]{y}$.
    \begin{align*}
        B &= \dfrac{x \cdot x^{\frac{1}{4}} \cdot y + x \cdot y \cdot y^{\frac{1}{4}}}{x^{\frac{1}{4}} + y^{\frac{1}{4}}} = \dfrac{xy\left(x^{\frac{1}{4}} + y^{\frac{1}{4}}\right)}{x^{\frac{1}{4}} + y^{\frac{1}{4}}} = xy.
    \end{align*}
\end{itemize}

\noindent \textbf{d) Tổ chức thực hiện:}
\begin{itemize}[leftmargin=0.6cm]
    \item \textbf{Giao nhiệm vụ:} Học sinh làm việc cá nhân 5 phút, sau đó thảo luận cặp đôi để thống nhất lời giải.
    \item \textbf{Thực hiện nhiệm vụ:} Giáo viên đi quanh lớp, hỗ trợ các bạn học sinh trung bình - yếu cách phân tích $x^{\frac{5}{4}} = x \cdot x^{\frac{1}{4}}$.
    \item \textbf{Báo cáo, thảo luận:} 2 học sinh lên bảng giải chi tiết. Cả lớp quan sát, nhận xét bước đặt nhân tử chung.
    \item \textbf{Kết luận, nhận định:} Giáo viên chốt lại phương pháp: Đưa các lũy thừa phân số về dạng tích của lũy thừa số mũ nguyên và căn thức để phát hiện nhân tử chung.
\end{itemize}

\subsubsection*{4. Hoạt động 4: Vận dụng (5 phút)}

\noindent \textbf{a) Mục tiêu:}
Vận dụng công thức lũy thừa số mũ hữu tỉ giải bài toán sinh học thực tiễn.

\noindent \textbf{b) Nội dung:}
\begin{stembox}{Bài toán Sinh học: Quy luật Kleiber về chuyển hóa năng lượng}
Theo quy luật Kleiber, tốc độ trao đổi chất cơ bản lúc nghỉ ngơi $R$ (tính bằng kcal/ngày) của một loài động vật có vú có liên hệ với khối lượng cơ thể $M$ (tính bằng kg) qua công thức lũy thừa:
$$R = 70 \cdot M^{\frac{3}{4}}$$
\begin{enumerate}[leftmargin=0.6cm]
    \item Một con linh dương có khối lượng $M = 16\text{ kg}$. Tính tốc độ trao đổi chất $R$ của con linh dương đó.
    \item Nếu khối lượng cơ thể tăng gấp $16$ lần thì tốc độ trao đổi chất $R$ tăng gấp bao nhiêu lần?
\end{enumerate}
\end{stembox}

\noindent \textbf{c) Sản phẩm:}
\begin{enumerate}[leftmargin=0.6cm]
    \item Với $M = 16\text{ kg}$:
    $$R = 70 \cdot 16^{\frac{3}{4}} = 70 \cdot (\sqrt[4]{16})^3 = 70 \cdot 2^3 = 70 \cdot 8 = 560\text{ (kcal/ngày)}.$$
    \item Khi khối lượng tăng gấp $16$ lần ($M' = 16M$):
    $$R' = 70 \cdot (16M)^{\frac{3}{4}} = 70 \cdot 16^{\frac{3}{4}} \cdot M^{\frac{3}{4}} = 8 \cdot \left(70 \cdot M^{\frac{3}{4}}\right) = 8R.$$
    Vậy tốc độ trao đổi chất tăng $8$ lần (chứ không phải tăng 16 lần như hàm tuyến tính).
\end{enumerate}

\noindent \textbf{d) Tổ chức thực hiện:}
\begin{itemize}[leftmargin=0.6cm]
    \item \textbf{Giao nhiệm vụ:} Giáo viên giao bài tập, gợi ý học sinh áp dụng công thức $(ab)^r = a^r b^r$.
    \item \textbf{Thực hiện nhiệm vụ:} Học sinh tính toán nhanh và ghi đáp án vào vở.
    \item \textbf{Báo cáo, nhận xét:} Học sinh chia sẻ câu trả lời; giáo viên chốt đáp số và dặn dò chuẩn bị cho Tiết 2.
\end{itemize}

\vspace{10pt}

\subsection*{\faClockO \ TIẾT 2 (TIẾT 56 THEO PPCT): LŨY THỪA VỚI SỐ MŨ THỰC, TÍNH CHẤT VÀ ỨNG DỤNG THỰC TIỄN}

\subsubsection*{1. Hoạt động 1: Mở đầu (Khởi động) (5 phút)}

\noindent \textbf{a) Mục tiêu:}
Kích thích trí tò mò khoa học về lũy thừa với số mũ vô tỉ (như $3^{\sqrt{2}}, 2^\pi$) và kết nối với kiến thức giới hạn dãy số đã học ở Chương V.

\noindent \textbf{b) Nội dung:}
Giáo viên đặt vấn đề:
Ta đã biết cách tính $3^2, 3^{\frac{1}{2}} = \sqrt{3}, 3^{\frac{1}{3}} = \sqrt[3]{3}$. Nhưng với số vô tỉ $\alpha = \sqrt{2} = 1{,}41421356\dots$, biểu thức $3^{\sqrt{2}}$ được định nghĩa như thế nào? Làm sao máy tính cầm tay tính được $3^{\sqrt{2}} \approx 4{,}728804388$?

\noindent \textbf{c) Sản phẩm:}
Học sinh dự đoán: Viết $\sqrt{2}$ dưới dạng dãy số thập phân hữu tỉ xấp xỉ:
$$r_1 = 1{,}4 = \dfrac{14}{10}; \quad r_2 = 1{,}41 = \dfrac{141}{100}; \quad r_3 = 1{,}414 = \dfrac{1414}{1000}; \quad \dots \quad \to \sqrt{2}$$
Khi đó dãy số các lũy thừa hữu tỉ:
$$3^{1{,}4} \approx 4{,}6555; \quad 3^{1{,}41} \approx 4{,}7069; \quad 3^{1{,}414} \approx 4{,}7276; \quad \dots$$
sẽ dần tiến tới một giới hạn xác định, đó chính là giá trị của $3^{\sqrt{2}}$!

\noindent \textbf{d) Tổ chức thực hiện:}
\begin{itemize}[leftmargin=0.6cm]
    \item \textbf{Giao nhiệm vụ:} Giáo viên trình chiếu bảng xấp xỉ, yêu cầu học sinh dùng MTCT bấm tính thử $3^{1{,}4}, 3^{1{,}41}, 3^{1{,}414}$ và so sánh với $3^{\sqrt{2}}$.
    \item \textbf{Thực hiện nhiệm vụ:} Học sinh bấm máy tính, quan sát hiện tượng giới hạn của dãy số.
    \item \textbf{Báo cáo, thảo luận:} Học sinh nêu nhận xét: Dãy các lũy thừa số mũ hữu tỉ tiến dần về một giá trị duy nhất.
    \item \textbf{Kết luận, nhận định:} Giáo viên chốt lại định nghĩa lũy thừa với số mũ thực thông qua giới hạn dãy số.
\end{itemize}

\subsubsection*{2. Hoạt động 2: Hình thành kiến thức (25 phút)}

\paragraph{2.1. Đơn vị kiến thức 1: Định nghĩa lũy thừa với số mũ thực (8 phút)}

\noindent \textbf{a) Mục tiêu:}
Hiểu được bản chất lũy thừa với số mũ thực của số thực dương là giới hạn của dãy số các lũy thừa số mũ hữu tỉ.

\noindent \textbf{b) Nội dung:}
\begin{pedabox}{Định nghĩa Lũy thừa với số mũ thực}
Cho $a$ là một số thực dương, $\alpha$ là một số vô tỉ.
Xét dãy số hữu tỉ $(r_n)$ thỏa mãn $\lim_{n \to +\infty} r_n = \alpha$.
Khi đó, dãy số $(a^{r_n})$ có giới hạn hữu hạn khi $n \to +\infty$.
Giới hạn đó không phụ thuộc vào việc chọn dãy số $(r_n)$ và được gọi là \textbf{lũy thừa của $a$ với số mũ $\alpha$}, ký hiệu là $a^\alpha$:
$$a^\alpha = \lim_{n \to +\infty} a^{r_n} \quad (\text{với } a > 0)$$
\end{pedabox}

\begin{scaffoldbox}{Tổng kết hệ thống điều kiện cơ số của lũy thừa $a^\alpha$}
\textbf{Bảng ghi nhớ bắt buộc:}
\begin{center}
\begin{tabular}{|l|c|c|}
    \hline
    \textbf{Loại số mũ $\alpha$} & \textbf{Điều kiện của cơ số $a$} & \textbf{Ví dụ} \\ \hline
    $\alpha \in \mathbb{N}^*$ (nguyên dương) & $a \in \mathbb{R}$ (tùy ý) & $(-3)^4 = 81, 0^5 = 0$ \\ \hline
    $\alpha \in \mathbb{Z} \setminus \mathbb{N}^*$ (nguyên âm hoặc $0$) & $a \neq 0$ & $2^0 = 1, (-5)^{-2} = \dfrac{1}{25}$ \\ \hline
    $\alpha \notin \mathbb{Z}$ (hữu tỉ không nguyên, vô tỉ, thực) & $a > 0$ (bắt buộc dương) & $4^{\frac{1}{2}} = 2, 5^{\sqrt{3}}, \pi^{-\sqrt{2}}$ \\ \hline
\end{tabular}
\end{center}
\end{scaffoldbox}

\noindent \textbf{c) Sản phẩm:}
Học sinh chỉ ra được các biểu thức có nghĩa và không có nghĩa:
\begin{itemize}
    \item $A = (x - 1)^0$ có nghĩa khi $x - 1 \neq 0 \iff x \neq 1$.
    \item $B = (2x - 4)^{-3}$ có nghĩa khi $2x - 4 \neq 0 \iff x \neq 2$.
    \item $C = (x - 3)^{\sqrt{2}}$ có nghĩa khi $x - 3 > 0 \iff x > 3$.
    \item $D = (x^2 + 1)^{-\frac{3}{5}}$ luôn có nghĩa với mọi $x \in \mathbb{R}$ vì $x^2 + 1 > 0 \ \forall x$.
\end{itemize}

\noindent \textbf{d) Tổ chức thực hiện:}
\begin{itemize}[leftmargin=0.6cm]
    \item \textbf{Giao nhiệm vụ:} Giáo viên trình bày định nghĩa và bảng điều kiện cơ số; yêu cầu học sinh tìm tập xác định của các biểu thức $A, B, C, D$.
    \item \textbf{Thực hiện nhiệm vụ:} Học sinh vận dụng bảng ghi nhớ để xác định điều kiện của $x$.
    \item \textbf{Báo cáo, thảo luận:} 4 học sinh lần lượt trả lời từng biểu thức. Giáo viên nhận xét, phân tích các lỗi học sinh dễ mắc.
    \item \textbf{Kết luận, nhận định:} Giáo viên nhấn mạnh: Cứ thấy số mũ không nguyên là phải đặt điều kiện biểu thức dưới cơ số $> 0$.
\end{itemize}

\paragraph{2.2. Đơn vị kiến thức 2: Tính chất của phép tính lũy thừa với số mũ thực (17 phút)}

\noindent \textbf{a) Mục tiêu:}
Nắm vững các tính chất của lũy thừa với số mũ thực; giải thích được tính chất so sánh hai lũy thừa cùng cơ số khi $a > 1$ và $0 < a < 1$.

\noindent \textbf{b) Nội dung:}
\begin{pedabox}{Các Tính chất của Phép tính Lũy thừa với số mũ thực}
Cho $a > 0, b > 0$ và $\alpha, \beta \in \mathbb{R}$. Khi đó:
\begin{enumerate}[leftmargin=0.6cm]
    \item $a^\alpha \cdot a^\beta = a^{\alpha + \beta}$
    \item $\dfrac{a^\alpha}{a^\beta} = a^{\alpha - \beta}$
    \item $(a^\alpha)^\beta = a^{\alpha \cdot \beta}$
    \item $(a \cdot b)^\alpha = a^\alpha \cdot b^\alpha$
    \item $\left(\dfrac{a}{b}\right)^\alpha = \dfrac{a^\alpha}{b^\alpha}$
    \item $a^0 = 1$; \quad $a^{-\alpha} = \dfrac{1}{a^\alpha}$; \quad $1^\alpha = 1$.
\end{enumerate}
\end{pedabox}

\begin{pedabox}{So sánh hai lũy thừa cùng cơ số}
Cho $a > 0$ và $\alpha, \beta \in \mathbb{R}$:
\begin{itemize}[leftmargin=0.6cm]
    \item Nếu $a > 1$ thì: $a^\alpha > a^\beta \iff \alpha > \beta$ (\textbf{Cùng chiều: cơ số lớn hơn 1, số mũ lớn hơn thì giá trị lớn hơn}).
    \item Nếu $0 < a < 1$ thì: $a^\alpha > a^\beta \iff \alpha < \beta$ (\textbf{Đổi chiều: cơ số trong khoảng $(0; 1)$, số mũ lớn hơn thì giá trị nhỏ hơn}).
\end{itemize}
\end{pedabox}

\noindent \textbf{c) Sản phẩm:}
Học sinh áp dụng giải các ví dụ:
\begin{itemize}
    \item So sánh $2^{\sqrt{3}}$ và $2^{\sqrt{2}}$: Vì cơ số $a = 2 > 1$ và $\sqrt{3} > \sqrt{2}$ nên $2^{\sqrt{3}} > 2^{\sqrt{2}}$.
    \item So sánh $\left(\dfrac{1}{3}\right)^{\pi}$ và $\left(\dfrac{1}{3}\right)^{3}$: Vì cơ số $a = \dfrac{1}{3} \in (0; 1)$ và $\pi \approx 3{,}14 > 3$ nên $\left(\dfrac{1}{3}\right)^\pi < \left(\dfrac{1}{3}\right)^3$.
    \item Rút gọn: $P = \dfrac{a^{\sqrt{3} + 1} \cdot a^{2 - \sqrt{3}}}{(a^{\sqrt{2} - 1})^{\sqrt{2} + 1}} = \dfrac{a^{(\sqrt{3} + 1) + (2 - \sqrt{3})}}{a^{(\sqrt{2} - 1)(\sqrt{2} + 1)}} = \dfrac{a^3}{a^{2 - 1}} = \dfrac{a^3}{a^1} = a^2$.
\end{itemize}

\noindent \textbf{d) Tổ chức thực hiện:}
\begin{itemize}[leftmargin=0.6cm]
    \item \textbf{Giao nhiệm vụ:} Giáo viên trình chiếu tính chất, yêu cầu học sinh làm việc theo nhóm 4 em hoàn thành 3 ví dụ trên.
    \item \textbf{Thực hiện nhiệm vụ:} Học sinh trao đổi trong nhóm, áp dụng tính chất lũy thừa và hằng đẳng thức $(A-B)(A+B) = A^2 - B^2$.
    \item \textbf{Báo cáo, thảo luận:} Mời 2 nhóm cử đại diện báo cáo cách làm.
    \item \textbf{Kết luận, nhận định:} Giáo viên nhấn mạnh tính chất so sánh: Phải luôn kiểm tra cơ số lớn hơn $1$ hay nhỏ hơn $1$ để giữ nguyên chiều hay đổi chiều bất đẳng thức.
\end{itemize}

\subsubsection*{3. Hoạt động 3: Luyện tập (10 phút)}

\noindent \textbf{a) Mục tiêu:}
Rèn luyện kỹ năng biến đổi rút gọn biểu thức chứa biến có số mũ thực và giải bài toán so sánh không dùng MTCT.

\noindent \textbf{b) Nội dung:}
Thực hiện Phiếu học tập số 2:
\begin{enumerate}[leftmargin=0.6cm]
    \item \textbf{Bài 1:} Cho $a > 0, b > 0$. Rút gọn biểu thức:
    $$K = \dfrac{a^{\frac{1}{3}} \cdot \sqrt{b} + b^{\frac{1}{3}} \cdot \sqrt{a}}{\sqrt[6]{a} + \sqrt[6]{b}}$$
    \item \textbf{Bài 2:} So sánh hai số sau mà không dùng MTCT: $A = (\sqrt{2} - 1)^{\sqrt{3}}$ và $B = (\sqrt{2} - 1)^{\sqrt{2}}$.
\end{enumerate}

\noindent \textbf{c) Sản phẩm:}
\begin{itemize}
    \item \textbf{Bài 1:}
    Quy đồng số mũ về mẫu chung là $6$:
    $a^{\frac{1}{3}} = a^{\frac{2}{6}} = \sqrt[6]{a^2}$; $\sqrt{b} = b^{\frac{1}{2}} = b^{\frac{3}{6}} = \sqrt[6]{b^3}$;
    $b^{\frac{1}{3}} = \sqrt[6]{b^2}$; $\sqrt{a} = \sqrt[6]{a^3}$.
    Tử số:
    $$\sqrt[6]{a^2} \cdot \sqrt[6]{b^3} + \sqrt[6]{b^2} \cdot \sqrt[6]{a^3} = \sqrt[6]{a^2 b^3} + \sqrt[6]{a^3 b^2} = \sqrt[6]{a^2 b^2}\left(\sqrt[6]{b} + \sqrt[6]{a}\right) = (ab)^{\frac{1}{3}} \left(\sqrt[6]{a} + \sqrt[6]{b}\right)$$
    Do đó:
    $$K = \dfrac{(ab)^{\frac{1}{3}} \left(\sqrt[6]{a} + \sqrt[6]{b}\right)}{\sqrt[6]{a} + \sqrt[6]{b}} = (ab)^{\frac{1}{3}} = \sqrt[3]{ab}.$$
    \item \textbf{Bài 2:}
    Xét cơ số: $\sqrt{2} \approx 1{,}414 \implies \sqrt{2} - 1 \approx 0{,}414 \in (0; 1)$.
    Xét hai số mũ: $\sqrt{3} > \sqrt{2} > 0$.
    Vì cơ số $0 < \sqrt{2} - 1 < 1$ và số mũ $\sqrt{3} > \sqrt{2}$, nên theo tính chất đổi chiều ta có:
    $$(\sqrt{2} - 1)^{\sqrt{3}} < (\sqrt{2} - 1)^{\sqrt{2}} \iff A < B.$$
\end{itemize}

\noindent \textbf{d) Tổ chức thực hiện:}
\begin{itemize}[leftmargin=0.6cm]
    \item \textbf{Giao nhiệm vụ:} Học sinh làm việc độc lập trong 6 phút.
    \item \textbf{Thực hiện nhiệm vụ:} Giáo viên quan sát, nhắc nhở học sinh cách quy đồng số mũ về phân số mẫu chung là $6$.
    \item \textbf{Báo cáo, thảo luận:} 2 học sinh trình bày bảng; lớp nhận xét và tự chấm điểm.
    \item \textbf{Kết luận, nhận định:} Giáo viên biểu dương học sinh có cách giải sáng tạo và ghi điểm đánh giá thường xuyên.
\end{itemize}

\subsubsection*{4. Hoạt động 4: Vận dụng (5 phút)}

\noindent \textbf{a) Mục tiêu:}
Vận dụng toán lũy thừa giải quyết bài toán tài chính thực tiễn (STEM) và mở rộng nhận thức về ứng dụng lũy thừa trong Trí tuệ nhân tạo (Mã 11.C5.1).

\noindent \textbf{b) Nội dung:}
\begin{stembox}{Bài toán STEM Kinh tế: Lãi kép liên tục}
Một người gửi tiết kiệm $100$ triệu đồng vào ngân hàng với lãi suất $r = 6\%$/năm ($r = 0{,}06$).
Nếu tiền lãi được tính theo thể thức \textbf{lãi kép liên tục} (lãi được cộng dồn vào gốc ở mọi thời điểm vô cùng bé), số tiền nhận được sau $t$ năm được tính theo công thức lũy thừa cơ số tự nhiên $e$:
$$A(t) = P \cdot e^{r \cdot t}$$
(trong đó $P = 100$ triệu đồng, $e \approx 2{,}71828$).
\begin{enumerate}[leftmargin=0.6cm]
    \item Tính số tiền cả gốc lẫn lãi người đó nhận được sau $5$ năm ($t = 5$).
    \item So sánh với thể thức lãi kép tính theo năm: $A_1 = P(1 + r)^t = 100 \cdot (1 + 0{,}06)^5$.
\end{enumerate}
\end{stembox}

\begin{aibox}{Tích hợp Năng lực AI (Theo Quyết định 2422/QĐ-BGDĐT - Mã 11.C5.1)}
\textbf{Ứng dụng toán lũy thừa trong Mạng nơ-ron nhân tạo (Neural Networks):}
\begin{itemize}[leftmargin=0.6cm]
    \item \textbf{Phép tính lũy thừa ma trận trọng số:} Khi dữ liệu truyền qua $L$ lớp ẩn trong mạng nơ-ron sâu (Deep Neural Network), tín hiệu đầu ra được tính bằng phép nhân lũy thừa ma trận trọng số $W^L$. Hiện tượng bùng nổ gradient (Exploding Gradient) hay biến mất gradient (Vanishing Gradient) xảy ra chính do tính chất hàm mũ $a^L$ ($a > 1 \implies a^L \to \infty$; $a < 1 \implies a^L \to 0$ khi $L$ lớn).
    \item \textbf{Hàm Softmax chuyển đổi xác suất:} Trợ lý AI (như ChatGPT, Gemini) phân loại từ ngữ tiếp theo dựa vào hàm mũ cơ số tự nhiên:
    $$P(y_i) = \dfrac{e^{z_i}}{\sum_{j} e^{z_j}}$$
    Nhờ tính chất lũy thừa $e^{z_i} > 0$ với mọi $z_i$, AI chuẩn hóa được điểm số thành xác suất phần trăm từ $0\%$ đến $100\%$.
\end{itemize}
\end{aibox}

\noindent \textbf{c) Sản phẩm:}
\begin{enumerate}[leftmargin=0.6cm]
    \item Với $t = 5$ năm:
    $$A(5) = 100 \cdot e^{0{,}06 \cdot 5} = 100 \cdot e^{0{,}3} \approx 100 \cdot 1{,}34986 = 134{,}986\text{ (triệu đồng)}.$$
    \item Với lãi kép hàng năm:
    $$A_1 = 100 \cdot (1{,}06)^5 \approx 100 \cdot 1{,}33823 = 133{,}823\text{ (triệu đồng)}.$$
    Ta thấy lãi kép liên tục đem lại lợi nhuận cao hơn khoảng $1{,}163$ triệu đồng sau 5 năm.
\end{enumerate}

\noindent \textbf{d) Tổ chức thực hiện:}
\begin{itemize}[leftmargin=0.6cm]
    \item \textbf{Giao nhiệm vụ:} Giáo viên giao bài toán STEM, hướng dẫn bấm số vô tỉ $e$ trên MTCT bằng tổ hợp phím \fbox{$\text{ALPHA}$} $\to$ \fbox{$\times 10^x$}.
    \item \textbf{Thực hiện nhiệm vụ:} Học sinh tính toán, thảo luận về ý nghĩa của công thức và lắng nghe phần giới thiệu ứng dụng AI.
    \item \textbf{Báo cáo, nhận xét:} Học sinh nêu kết quả, giáo viên tổng kết bài học, giao bài tập về nhà trong SGK và chuẩn bị cho Bài 19: Lôgarit.
\end{itemize}

\vspace{10pt}

\section*{IV. PHỤ LỤC HỌC LIỆU BỔ TRỢ (BẮT BUỘC)}

\subsection*{1. Phiếu học tập số 1 (Tiết 1): Phân tầng bậc thang Lũy thừa số mũ nguyên và hữu tỉ}

\begin{pedabox}{PHIẾU HỌC TẬP SỐ 1: KHÁM PHÁ LŨY THỪA SỐ MŨ HỮU TỈ}
\textbf{Họ và tên học sinh:} \dotfill \ \textbf{Lớp:} \dotfill \ \textbf{Nhóm:} \dotfill

\noindent \textbf{CẤP ĐỘ 1: NHẬN BIẾT (Khởi động cơ bản)}
\begin{enumerate}[leftmargin=0.6cm]
    \item Tính giá trị các lũy thừa số mũ nguyên sau:
    $$a) \ (-5)^0 = \dots\dots\dots; \qquad b) \ 4^{-3} = \dots\dots\dots; \qquad c) \ \left(\dfrac{2}{5}\right)^{-2} = \dots\dots\dots$$
    \item Đưa các căn thức sau về dạng lũy thừa số mũ hữu tỉ ($a > 0$):
    $$a) \ \sqrt[3]{a^2} = \dots\dots\dots; \qquad b) \ \sqrt[5]{a^7} = \dots\dots\dots; \qquad c) \ \dfrac{1}{\sqrt[4]{a^3}} = \dots\dots\dots$$
\end{enumerate}

\noindent \textbf{CẤP ĐỘ 2: THÔNG HIỂU (Vận dụng quy tắc \& Thao tác MTCT)}
\begin{enumerate}[leftmargin=0.6cm]
    \item Tính giá trị biểu thức số (tính tay rồi dùng MTCT Casio fx-580VN X kiểm tra đối chiếu):
    $$P = 8^{\frac{4}{3}} + 81^{\frac{3}{4}} - \left(\dfrac{1}{16}\right)^{-0{,}75}$$
    \item Rút gọn biểu thức chứa biến với $x > 0$:
    $$Q = \dfrac{x^{\frac{4}{3}} \cdot \sqrt{x}}{\sqrt[3]{x^2} \cdot x^{\frac{1}{6}}}$$
\end{enumerate}

\noindent \textbf{CẤP ĐỘ 3: VẬN DỤNG (Thử thách phát triển tư duy)}
\begin{enumerate}[leftmargin=0.6cm]
    \item Cho biểu thức $M = \sqrt{x \sqrt[3]{x \sqrt[4]{x}}}$ với $x > 0$. Hãy viết biểu thức $M$ dưới dạng lũy thừa $x^\alpha$. Tính giá trị của $\alpha$.
\end{enumerate}
\end{pedabox}

\subsection*{2. Phiếu học tập số 2 (Tiết 2): Rút gọn biểu thức số mũ thực \& Ứng dụng thực tiễn}

\begin{pedabox}{PHIẾU HỌC TẬP SỐ 2: CHINH PHỤC LŨY THỪA SỐ MŨ THỰC}
\textbf{Họ và tên học sinh:} \dotfill \ \textbf{Lớp:} \dotfill \ \textbf{Nhóm:} \dotfill

\noindent \textbf{CẤP ĐỘ 1: NHẬN BIẾT \& THÔNG HIỂU (Điều kiện cơ số \& So sánh)}
\begin{enumerate}[leftmargin=0.6cm]
    \item Tìm điều kiện xác định của các biểu thức sau:
    $$a) \ A(x) = (2x - 6)^{\sqrt{5}}; \qquad b) \ B(x) = (x^2 - 4)^{-2}; \qquad c) \ C(x) = (x^2 - 3x + 2)^{\frac{1}{3}}$$
    \item Điền dấu thích hợp ($>, <, =$) vào ô trống không dùng máy tính cầm tay:
    $$a) \ 3^{\sqrt{5}} \ \dots\dots \ 3^{\sqrt{3}}; \qquad b) \ (0{,}5)^{\pi} \ \dots\dots \ (0{,}5)^{3}; \qquad c) \ (\sqrt{3} - 1)^{-\sqrt{2}} \ \dots\dots \ (\sqrt{3} - 1)^{-\sqrt{3}}$$
\end{enumerate}

\noindent \textbf{CẤP ĐỘ 2: VẬN DỤNG (Rút gọn biểu thức phức tạp)}
\begin{enumerate}[leftmargin=0.6cm]
    \item Rút gọn biểu thức với $a > 0, b > 0$:
    $$E = \left(a^{\sqrt{2}} - b^{\sqrt{2}}\right)^2 + 4(ab)^{\frac{\sqrt{2}}{2}} - (a^{\sqrt{2}} + b^{\sqrt{2}})^2$$
    \item Cho biết $a^{\frac{1}{2}} + a^{-\frac{1}{2}} = 3$. Hãy tính giá trị của các biểu thức sau:
    $$a) \ a + a^{-1}; \qquad b) \ a^2 + a^{-2}$$
\end{enumerate}
\end{pedabox}

\subsection*{3. Bảng Rubric đánh giá năng lực học sinh theo 4 mức độ}

\begin{center}
\begin{tabularx}{\textwidth}{|p{2.8cm}|X|X|X|X|}
    \hline
    \textbf{Tiêu chí} & \textbf{Mức 1: Chưa đạt} & \textbf{Mức 2: Đạt} & \textbf{Mức 3: Khá} & \textbf{Mức 4: Tốt} \\ \hline
    \textbf{1. Hiểu điều kiện cơ số và tính chất lũy thừa} & Hay nhầm lẫn điều kiện cơ số; cho rằng $(-2)^{\frac{1}{2}}$ có nghĩa. & Nắm được cơ số phải dương khi số mũ không nguyên; nhớ các công thức cơ bản. & Phân biệt chính xác điều kiện cho từng loại số mũ; áp dụng đúng tính chất so sánh. & Giải thích sâu sắc lý do tại sao cơ số phải dương; liên hệ tốt với tập xác định hàm số. \\ \hline
    \textbf{2. Kỹ năng tính toán \& Rút gọn biểu thức} & Lúng túng khi biến đổi căn thức về số mũ hữu tỉ; tính sai số mũ âm. & Biến đổi được các biểu thức đơn giản; tính đúng giá trị số cơ bản. & Rút gọn thành thạo biểu thức chứa biến nhiều tầng căn; quy đồng số mũ chuẩn xác. & Biến đổi nhanh, linh hoạt; phát hiện nhân tử chung và hằng đẳng thức lũy thừa sáng tạo. \\ \hline
    \textbf{3. Năng lực số \& Sử dụng MTCT (Mã 5.1NC1a)} & Chưa biết bấm phím lũy thừa phân số; không giải thích được lỗi \texttt{Math ERROR}. & Bấm được lũy thừa cơ bản theo hướng dẫn của giáo viên và bạn bè. & Thao tác bấm máy nhanh, chính xác; dùng MTCT kiểm tra chéo kết quả rút gọn. & Điêu luyện trong sử dụng MTCT; giải thích được cơ chế xử lý xấp xỉ số mũ vô tỉ của máy. \\ \hline
    \textbf{4. Năng lực mô hình hóa \& Năng lực AI (Mã 11.C5.1)} & Chưa liên hệ được toán lũy thừa với các bài toán thực tiễn đời sống. & Thay số và tính được bài toán lãi kép cơ bản khi có công thức cho sẵn. & Tự thiết lập được mô hình tăng trưởng lũy thừa; giải thích được ý nghĩa các đại lượng. & Hiểu rõ cơ chế lũy thừa trong thuật toán AI, mạng nơ-ron; liên hệ thực tiễn xuất sắc. \\ \hline
\end{tabularx}
\end{center}

\end{document}
